% Generated draft: requires mathematical and source review.
\documentclass[11pt]{article}
\usepackage[margin=25mm]{geometry}
\usepackage{fontspec}
\setmainfont{Latin Modern Roman}
\usepackage{amsmath,amssymb,mathtools,unicode-math}
\setmathfont{Latin Modern Math}
\usepackage{enumitem,needspace,xurl,hyperref}
\setlength{\parindent}{0pt}
\setlength{\parskip}{5pt}
\setlength{\emergencystretch}{3em}
\newcommand{\N}{\mathbb{N}}
\newcommand{\Z}{\mathbb{Z}}
\newcommand{\Q}{\mathbb{Q}}
\newcommand{\R}{\mathbb{R}}
\newcommand{\C}{\mathbb{C}}
\newcommand{\F}{\mathbb{F}}
\newcommand{\A}{\mathbb{A}}
\newcommand{\PP}{\mathbb{P}}
\newcommand{\E}{\mathbb{E}}
\newcommand{\eps}{\varepsilon}
\newcommand{\dd}{\,\mathrm{d}}
\begin{document}
\section*{UnsolvedMath ID: 30005701}
\textbf{Source problem code:} catalogue entry
\subsection{Definitions and mathematical statement}
For a finitely generated group \(G\), a real character is a homomorphism \(\chi:G\to(\R,+)\). Its character sphere is
\[
S(G)=\bigl(\operatorname{Hom}(G,\R)\setminus\{0\}\bigr)/{\sim},
\qquad \chi\sim\chi'\iff \chi'=a\chi\text{ for some }a>0.
\]
Write \([\chi]\) for the resulting positive ray.

One standard interpretation of the Bieri--Neumann--Strebel--Renz invariants is the homotopical family \(\Sigma^n(G)\subseteq S(G)\). Here \(n\geq1\), and \(G\) must have type \(F_n\): there must exist an Eilenberg--Mac Lane space \(K(G,1)\) with finite \(n\)-skeleton. Choose such a space, let \(X\) be the \(n\)-skeleton of its universal cover, and choose a continuous height function \(h_\chi:X\to\R\) satisfying
\[
h_\chi(gx)=h_\chi(x)+\chi(g).
\]
Put \(X_{\chi,t}=h_\chi^{-1}([t,\infty))\). The standard controlled-connectivity definition is
\[
[\chi]\in\Sigma^n(G)
\iff
\begin{gathered}
\text{for every }t\in\R\text{ there exists }s\leq t\text{ such that, for every }0\leq j\leq n-1,\\
\text{every continuous }f:S^j\to X_{\chi,t}\text{ extends to a continuous }\\
\bar f:B^{j+1}\to X_{\chi,s}.
\end{gathered}
\]
Here \(S^j=\partial B^{j+1}\), and an extension agrees with \(f\) on that boundary. The threshold \(s\) must work for every map in the indicated dimensions. Homological invariants \(\Sigma^n(G;R)\), with a specified coefficient ring \(R\) and corresponding finiteness hypotheses, are another convention; the supplied record does not select between these conventions.

\textbf{Source problem.} Determine the Bieri--Neumann--Strebel--Renz invariant of every algebraic Bieri--Strebel group, extending the computation for the group denoted \(F_\tau\) in the source. Under the homotopical, all-degree interpretation, this would mean giving an explicit description of
\[
\Sigma^n(G)\subseteq S(G)
\qquad\text{for every intended algebraic Bieri--Strebel group }G
\text{ and every }n\geq1\text{ for which }G\text{ has type }F_n.
\]
This displayed interpretation is conditional: the supplied record does not define the class of algebraic Bieri--Strebel groups, its admissible parameters, or any interval, slope, breakpoint, or endpoint conditions. It also does not specify the degrees, a homotopical or homological convention, or homological coefficients. Consequently, the supplied evidence does not determine a fully specified quantified problem beyond the source target stated above. The notation \(F_\tau\) is retained without assigning an unsupported definition to the group or a value to \(\tau\).

The supplied account credits Molyneux, Nucinkis and Santos Rego with computing the invariant of \(F_\tau\) using a method of Bieri, Geoghegan and Kochloukova. Its dated assessment also reports computations for further subclasses. These are supplied background claims, not computations or literature conclusions established by this draft.
\subsection{Short English statement}
The problem asks for a general description of the BNSR invariants of algebraic Bieri--Strebel groups, extending the reported golden-mean example \(F_\tau\). The source omits the definition of the group class and the precise invariant convention, so the homotopical formulation above is conditional. The supplied open status and reported special-case computations require independent review.
\subsection{Sources}
\begin{itemize}
\item Local UnsolvedMath registry: supplied complete source record, permanent ID \(30005701\), exact source problem code \texttt{catalogue entry}, title \emph{BNSR Invariants of Algebraic Bieri--Strebel Groups}, in the Oberwolfach Reports Open Problems collection.
\item Supplied report citation: \emph{Mini-Workshop: Homological aspects of TDLC-groups (2024)}, Oberwolfach Reports, DOI \texttt{10.4171/owr/2023/55}, pp.\ 3103--3138. \url{https://doi.org/10.4171/owr/2023/55}. The registry separately gives proposed year \(2023\) and an \texttt{owr-year:2023} tag; the relationship between these dates and the citation's \(2024\) label is unresolved here.
\item Supplied attribution: Lewis Molyneux, \emph{Finiteness Properties of Algebraic Bieri--Strebel Groups}, joint work with Brita Nucinkis and Yuri Santos Rego. The excerpt attributes the \(F_\tau\) calculation to Molyneux, Nucinkis and Santos Rego, refers to it as reference [7], and credits the method to Bieri, Geoghegan and Kochloukova. The excerpt does not supply the full bibliographic entry for [7].
\item Supplied literature source, labelled ``Oberwolfach report'': \url{https://ems.press/content/serial-article-files/48166}.
\item Supplied literature source, labelled ``Molyneux--Nucinkis--Santos Rego'': \url{https://nyjm.albany.edu/j/2024/30-23.html}.
\item Supplied literature source, labelled ``Molyneux'': \url{https://arxiv.org/abs/2602.08748}.
\item Literature assessment dated exactly \texttt{2026-08-22}: the supplied curation labels the general problem open and reports computations for the golden-mean case and further subclasses. This is supplied dated curation, not an independently verified assessment. The background explicitly calls it dated literature triage derived from the supplied OWR curation and requests independent verification of time-sensitive status and cited claims.
\item The registry labels the statement \texttt{corrected\_verified} and records a claim that it was checked against the official report. This draft preserves that provenance claim without independently confirming it or visiting any cited source.
\end{itemize}
\end{document}
